\documentclass[12pt,a4paper]{article}

% --- PAQUETES DE CONFIGURACIÓN ---
\usepackage[utf8]{inputenc}
\usepackage[spanish,es-tabla]{babel}
\usepackage[margin=2.5cm]{geometry}
\usepackage{graphicx}
\usepackage{booktabs}
\usepackage{array}
\usepackage{tabularx}
\usepackage{enumitem}
\usepackage{xcolor}
\usepackage{titlesec}
\usepackage{fancyhdr}
\usepackage{hyperref}
\usepackage{float} % Permite fijar tablas y figuras exactamente donde se definen usando [H]

% Configuración de Hipervínculos
\hypersetup{
    colorlinks=true,
    linkcolor=blue,
    filecolor=magenta,      
    urlcolor=blue,
    pdftitle={Prueba Práctica Unidad IV - Ingeniería de Requisitos},
    pdfpagemode=FullScreen,
}

% Estilo de Encabezados y Pies de Página
\pagestyle{fancy}
\fancyhf{}
\rhead{\footnotesize Ingeniería de Requisitos (ISR-401) | Prueba Práctica U-IV}
\lhead{\footnotesize UTEQ - Ingeniería de Software}
\cfoot{\thepage}

\setlength{\parindent}{0pt}
\setlength{\parskip}{6pt}

\begin{document}

% ==============================================================================
% PORTADA INSTITUCIONAL UTEQ
% ==============================================================================
\begin{titlepage}
    \centering
    {\bfseries\Large UNIVERSIDAD TÉCNICA ESTATAL DE QUEVEDO \par}
    \vspace{0.2cm}
    {\bfseries\large FACULTAD DE CIENCIAS DE LA INGENIERÍA \par}
    {\bfseries\large CARRERA DE INGENIERÍA DE SOFTWARE \par}
    \vspace{0.4cm}
    
    \rule{\linewidth}{0.5mm} \\[0.3cm]
    {\huge\bfseries PRUEBA PRÁCTICA UNIDAD IV \par}
    \vspace{0.2cm}
    {\large\bfseries Validación, Gestión, Herramientas y Estándares de Requisitos \par}
    \rule{\linewidth}{0.5mm} \\[0.6cm]
    
    \begin{minipage}{0.48\textwidth}
        \begin{flushleft} \large
            \textbf{Asignatura:} Ingeniería de Requisitos (ISR-401)\\
            \textbf{Estudiante:} Alvia Villegas Erick Adalberto\\
            \textbf{Docente:} Ing. Gleiston Guerrero, Mg.\\
            \textbf{Modalidad:} Individual, en clase
        \end{flushleft}
    \end{minipage}
    \begin{minipage}{0.48\textwidth}
        \begin{flushright} \large
            \textbf{Fecha:} 11/08/2026\\
            \textbf{Duración:} 120 minutos\\
            \textbf{Caso:} Gestión de Pedidos\\
            \textbf{Ponderación:} 100 pts (30+70)
        \end{flushright}
    \end{minipage}
    
    \vspace{0.5cm}
    
    % URL DEL REPOSITORIO
    \begin{colorbox}{gray!15}
        \centering
        \vspace{0.15cm}
        \textbf{URL del Repositorio Git Público (Criterio de Piso):}\\
        \url{https://github.com/Erick-Alvia/SGCV-IA-Practica-U4.git}
        \vspace{0.15cm}
    \end{colorbox}
    
    \vspace{0.5cm}
    
    % SECCIÓN AMPLIADA PARA CAPTURAS SGA DE EVIDENCIA
    \renewcommand{\arraystretch}{1.2}
    \begin{tabularx}{\linewidth}{|X|X|}
        \hline
        \hfil \textbf{Captura Resumen del Cuestionario SGA} & \hfil \textbf{Captura de Evaluación (Intento SGA)} \tabularnewline
        \hline
        \vspace{0.1cm}
        \centering
        \includegraphics[width=0.95\linewidth, height=6.2cm, keepaspectratio]{Resumen_Cuestionario.png}
        \vspace{0.1cm}
        & 
        \vspace{0.1cm}
        \centering
        \includegraphics[width=0.95\linewidth, height=6.2cm, keepaspectratio]{Revision_Intento.png}
        \vspace{0.1cm}
        \tabularnewline
        \hline
    \end{tabularx}

    \vfill
    {\small Quevedo - Los Ríos - Ecuador \par}
    {\small 2026 \par}
\end{titlepage}

\newpage

% ==============================================================================
% DESARROLLO DE LAS ACTIVIDADES PRÁCTICAS (P1 - P10)
% ==============================================================================

\section*{Caso Práctico: Sistema de Gestión de Pedidos}

A continuación se presenta el desarrollo de las diez actividades prácticas derivadas del dominio del \textbf{Sistema de Gestión de Pedidos}, manteniendo total consistencia entre los modelos funcional, de datos, de comportamiento y de trazabilidad.

---

\subsection*{P1. Modelo de datos - Diagrama de clases UML (8 pts)}
\textbf{Instrucción:} Construya el diagrama de clases UML del dominio con las clases \texttt{Cliente}, \texttt{Pedido}, \texttt{Linea\_Pedido} y \texttt{Producto}, indicando atributos, operaciones, multiplicidades e identificadores subrayados.

\begin{figure}[H]
    \centering
    \includegraphics[width=0.8\linewidth]{Class_Diagram.png}
    \caption{Diagrama de Clases UML — Dominio del Sistema de Pedidos.}
    \label{fig:p1_diagrama_clases}
\end{figure}

---
\newpage

\subsection*{P2. Modelo funcional - Diagrama de actividades UML (8 pts)}
\textbf{Instrucción:} Modele el proceso "Registrar pedido" utilizando carriles de responsabilidad (swimlanes), nodos de decisión, flujos alternativos, convergencia y nodo final.

\begin{figure}[H]
    \centering
    \includegraphics[width=0.8\linewidth]{Activity_DIagram.png}
    \caption{Diagrama de Actividades UML — Proceso Registrar Pedido.}
    \label{fig:p2_diagrama_actividades}
\end{figure}


\subsection*{P3. Modelo de comportamiento - Máquina de estados UML (8 pts)}
\textbf{Instrucción:} Elaborar la máquina de estados del ciclo de vida de un \texttt{Pedido}. Utilizar un superestado para factorizar la acción de anulación.

\begin{figure}[H]
    \centering
    \includegraphics[width=0.8\linewidth]{Machine_State.png}
    \caption{Máquina de Estados UML — Ciclo de Vida de un Pedido.}
    \label{fig:p3_maquina_estados}
\end{figure}

\newpage

\subsection*{P4. Consistencia entre las tres perspectivas (7 pts)}
\textbf{Instrucción:} Presentar la tabla de verificación cruzada entre P1, P2 y P3, identificando y corrigiendo inconsistencias.

\begin{table}[H]
\centering
\small
\begin{tabularx}{\textwidth}{|X|X|X|}
\hline
\textbf{Evento (P3)} & \textbf{Función / Acción (P2)} & \textbf{Entidad / Atributo modificado (P1)} \\
\hline
ev: verificarStock & Verifica stock / Revisa Stock & Producto.stock \\
\hline
ev: confirmar & Confirmar Pedido & Pedido.total, Pedido.estado \\
\hline
ev: notificarFaltante & Notificar faltante & Pedido.estado \\
\hline
ev: despachar / entregar / anular & Acciones de seguimiento & Pedido.estado \\
\hline
\end{tabularx}
\end{table}

\textbf{Defecto Detectado:} El evento \texttt{ev:notificarFaltante} en P3 no poseía un atributo de control explícito de estado en la clase \texttt{Pedido} dentro de P1 para registrar dicha notificación.  
\textbf{Corrección:} Se añade formalmente el atributo implícito de control de estados en la clase \texttt{Pedido}.


\subsection*{P5. Especificación de requisitos con esquema de atributos (10 pts)}
\textbf{Instrucción:} Redactar 2 requisitos funcionales y 2 no funcionales (métrica/umbral ISO/IEC 25010:2023) con su esquema de atributos completo.

\begin{table}[H]
\centering
\small
\begin{tabularx}{\textwidth}{|l|l|X|}
\hline
\textbf{ID} & \textbf{Tipo} & \textbf{Descripción / Métrica y Umbral (ISO 25010)} \\
\hline
\textbf{RF-01} & Funcional & El sistema debe permitir registrar un pedido verificando la disponibilidad de stock de cada producto seleccionado. \\
\hline
\textbf{RF-02} & Funcional & El sistema debe permitir anular un pedido en estado Confirmado o Enviado antes de ser entregado. \\
\hline
\textbf{RNF-01} & No Funcional & \textbf{ISO 25010: Eficiencia de Desempeño.} El tiempo de respuesta al registrar un pedido debe ser $<$ 3 segundos bajo una carga normal de 100 usuarios concurrentes. \\
\hline
\textbf{RNF-02} & No Funcional & \textbf{ISO 25010: Seguridad (Protección de datos).} Los datos personales del cliente (cédula, correo) deben cifrarse mediante AES-256 en reposo y tránsito. \\
\hline
\end{tabularx}
\end{table}

\newpage

\subsection*{P6. Priorización MoSCoW (5 pts)}
\textbf{Instrucción:} Clasificar los requisitos de P5 utilizando la técnica MoSCoW, justificando según valor vs. viabilidad.

\begin{table}[H]
\centering
\small
\begin{tabularx}{\textwidth}{|l|c|X|}
\hline
\textbf{Req.} & \textbf{Categoría} & \textbf{Justificación breve (Valor vs. Viabilidad)} \\
\hline
\textbf{RF-01} & \textbf{Must} & Funcionalidad operativa core imprescindible para el funcionamiento del negocio. \\
\hline
\textbf{RNF-02} & \textbf{Must} & Garantiza el cumplimiento normativo legal de protección de datos personales. \\
\hline
\textbf{RNF-01} & \textbf{Should} & Aporta alta satisfacción y usabilidad al usuario; alcanzable técnicamente. \\
\hline
\textbf{RF-02} & \textbf{Could} & Aporta valor de gestión pero puede procesarse manualmente en la primera versión. \\
\hline
\end{tabularx}
\end{table}

---

\subsection*{P7. Validación por inspección (lista de comprobación 29148) (8 pts)}
\textbf{Instrucción:} Registrar defectos aplicando las propiedades del estándar IEEE 29148 y presentar el retrabajo posterior.

\subsubsection*{1. Registro de Defectos Detectados (Sin corregir)}
\begin{itemize}[leftmargin=*]
    \item \textbf{Defecto RF-01:} Ambigüedad al no detallar explícitamente la acción del sistema ante una falta parcial de stock durante el registro.
    \item \textbf{Defecto RNF-01:} Falta de verificabilidad precisa al omitir las condiciones de infraestructura del servidor de prueba.
\end{itemize}

\subsubsection*{2. Retrabajo (Requisitos Corregidos)}
\begin{itemize}[leftmargin=*]
    \item \textbf{RF-01 (Corregido):} El sistema debe registrar un pedido validando el stock; si la cantidad solicitada excede el stock disponible, notificará el faltante y bloqueará la confirmación del pedido.
    \item \textbf{RNF-01 (Corregido):} El sistema debe procesar la transacción de registro de pedido en un tiempo $<$ 3s en el servidor de pruebas estándar.
\end{itemize}

\newpage

\subsection*{P8. Pruebas de aceptación trazadas (6 pts)}
\textbf{Instrucción:} Definir una prueba de aceptación por requisito indicando tipo, entradas y criterios de éxito.

\begin{table}[H]
\centering
\small
\begin{tabularx}{\textwidth}{|l|l|l|X|X|}
\hline
\textbf{ID Prueba} & \textbf{Req} & \textbf{Tipo} & \textbf{Entradas de Prueba} & \textbf{Criterio de Éxito} \\
\hline
\textbf{TC-01} & RF-01 & Funcional & Cliente válido, Producto A (stock=5), Cantidad=2. & Pedido registrado en estado Confirmado y stock actualizado a 3. \\
\hline
\textbf{TC-02} & RF-02 & Funcional & ID de pedido en estado Enviado, acción de anulación. & Estado del pedido cambia a Cancelado y se libera el stock. \\
\hline
\textbf{TC-03} & RNF-01 & No Funcional & 100 solicitudes simultáneas mediante JMeter. & 95\% de las respuestas se reciben en un tiempo $<$ 3 segundos. \\
\hline
\textbf{TC-04} & RNF-02 & No Funcional & Inspección de BD y tráfico de red durante registro. & Cédula y correo cifrados en BD (AES-256) y TLS 1.3 activo. \\
\hline
\end{tabularx}
\end{table}

---

\subsection*{P9. Matriz de trazabilidad (6 pts)}
\textbf{Instrucción:} Construir la matriz cruzando requisitos con fuentes, modelos UML, pruebas y dependencias horizontales.

\begin{table}[H]
\centering
\small
\begin{tabularx}{\textwidth}{|l|X|X|l|l|c|}
\hline
\textbf{Requisito} & \textbf{Fuente (Atrás)} & \textbf{Modelos UML (Adelante)} & \textbf{Prueba} & \textbf{Dep. Horiz.} & \textbf{¿Huérfano?} \\
\hline
\textbf{RF-01} & Cliente / Tienda & P1 (Pedido, Producto), P2, P3 & TC-01 & RNF-01 & No \\
\hline
\textbf{RF-02} & Tienda / Políticas & P1 (Pedido), P3 (Superestado) & TC-02 & RF-01 & No \\
\hline
\textbf{RNF-01} & Calidad / SLA & P2 (Proceso Registrar) & TC-03 & Ninguna & No \\
\hline
\textbf{RNF-02} & Seguridad / Ley & P1 (Cliente) & TC-04 & Ninguna & No \\
\hline
\end{tabularx}
\end{table}

\newpage

\subsection*{P10. Gestión del cambio y línea base (4 pts)}
\textbf{Instrucción:} Simular una Solicitud de Cambio (RFC), analizar su impacto, registrar la decisión del CCB y declarar la línea base.

\subsubsection*{1. Solicitud de Cambio (RFC-01)}
\textbf{Descripción:} Modificar \texttt{RF-02} para permitir la anulación de pedidos que ya estén en estado Entregado dentro de las primeras 24 horas posteriores a la entrega.

\subsubsection*{2. Análisis de Impacto}
Impacta a \textbf{P1} (requiere añadir el atributo \texttt{fechaEntrega}), a \textbf{P3} (requiere una nueva transición entre \texttt{Entregado} y \texttt{Cancelado}), al caso de prueba \textbf{TC-02} y a las políticas operativas del negocio.

\subsubsection*{3. Decisión del CCB (Change Control Board)}
\textbf{Decisión:} \textbf{Diferida / Rechazada para la versión actual.} Se posterga su implementación para la versión 2.0 debido al nivel de complejidad técnica y operativa que introduce sobre el proceso logístico desplegado.

\subsubsection*{4. Declaración de Línea Base (Baseline v1.0)}
Se declara y congela formalmente la configuración inicial compuesta por:
\begin{itemize}
    \item \textbf{Especificación ERS:} v1.0 (Requisitos RF-01, RF-02, RNF-01, RNF-02)
    \item \textbf{Modelos UML:} v1.0 (P1 Clases, P2 Actividades, P3 Estados)
    \item \textbf{Matriz de Trazabilidad:} v1.0
    \item \textbf{Pruebas de Aceptación:} v1.0 (TC-01 a TC-04)
\end{itemize}

\end{document}
